% ID: SEL-D0069-OA-E02
% ESTADO: validada. S069.
\textbf{a)} El dominio es $\mathbb R$. El denominador $x+1$ es positivo en el tramo donde se usa. Las ramas son suaves en sus interiores. En $0$, el valor y ambos límites son $0$, así que hay continuidad. Para la derivada en la unión, el cociente incremental es $x+1$ a la izquierda y $1/(x+1)$ a la derecha. Los dos límites valen $1$, así que $f'(0)=1$. En los interiores las derivadas son $2x+1$ y $1/(x+1)^2$, respectivamente.

\textbf{b)} A la derecha, $x/(x+1)=1-1/(x+1)$ tiende a $1$, por lo que $\boxed{y=1}$ es asíntota horizontal por la derecha. A la izquierda, $x^2+x$ tiende a $+\infty$ y no hay asíntota horizontal.

\textbf{c)} No hay asíntota vertical porque la función es continua en todo $\mathbb R$. En particular, en $-1$ se aplica la rama polinómica.
